	\renewcommand{\labelenumi}{\alph{enumi}.)}
	\renewcommand{\labelenumii}{\arabic*.}
{\scriptsize 	\begin{tabularx}{\textwidth}{lX}
		\textit{Vorausgesetzte Techniken}:& Zuki, Age-Uke, Yoko-Uke, Harai-Otoshi-Uke, Soto-Uke, Uraken-Uchi, Empi-Age-Uchi, Mae-Geri, \mbox{\textit{(Chisai-no-)}Mawashi-Geri}
	\end{tabularx}\\}
	\begin{enumerate}[nosep]
		\item \textit{\textbf{Kihon-Ido}}:
		\begin{enumerate}[nosep]
			\item Harai-Otoshi-Uke/Gyaku-Zuki Ch\={u}dan\tabto{280pt}(Zenkutsu-Dachi)
			\item Oi-Zuki J\={o}dan/Gyaku-Zuki Ch\={u}dan\tabto{280pt}(Zenkutsu-Dachi)
			\item Soto-Uke Ch\={u}dan/Gyaku-Zuki Ch\={u}dan\tabto{280pt}(Zenkutsu-Dachi)
			\item Kansetsu-Geri Gedan\tabto{280pt}(Sanchin-Dachi)
			\item Age-Uke/Gyaku-Zuki J\={o}dan\tabto{280pt}(Sanchin-Dachi)
			\item Yoko-Uke/Gyaku-Zuki Ch\={u}dan\tabto{280pt}(Sanchin-Dachi)
			\item Harai-Otoshi-Uke/Gyaku-Zuki Gedan\tabto{280pt}(Shiko-Dachi 45°)
		\end{enumerate}
		\item \textit{\textbf{Kata}}:
			\begin{enumerate}[nosep]
			\item Taikyoku (welche Taikyoku geprüft wird, entscheidet der Prüfer)
			\item Gekisai-Dai-Ichi
			\end{enumerate}
		\item \textit{\textbf{Partnerformen}}:
		\begin{enumerate}[nosep]
			\item Kihon-Ippon Kumite:\tabto{280pt}J\={o}dan\&Ch\={u}dan (ggf. Gedan)
		\end{enumerate}
	\end{enumerate}